\Askhsh{24758}{4}
\begin{ylh}{1}
    \item 2.1 Η έννοια της παραγώγου
    \item 2.2 Παραγωγίσιμες συναρτήσεις - Παράγωγος συνάρτηση
    \item 2.3 Κανόνες παραγώγισης
    \item 2.7 Τοπικά ακρότατα συνάρτησης
    \item 2.8 Κυρτότητα - Σημεία καμπής συνάρτησης
    \item 3.4 Ορισμένο ολοκλήρωμα.
\end{ylh}
Έστω συνάρτηση $f: \mathbb{R} \to \mathbb{R}$ παραγωγίσιμη με συνεχή παράγωγο, και η συνάρτηση $g(x)=(x^2-1)f(x)$ για την οποία ισχύει $g(x) \ge 0$ για κάθε $x \in \mathbb{R}$. Να αποδείξετε ότι:
\begin{alist}
\item η $g$ παρουσιάζει ελάχιστο για $x=1$ και για $x=-1$ και στη συνέχεια ότι $f(1)=f(-1)=0$. \monades{6}
\item $f'(1) \ge 0$ και $f'(-1) \le 0$. \monades{8}
\item η $f$ δεν είναι κοίλη. \monades{5}
\item $\int_{-1}^{1} (x^3-3x) f'(x)dx \le 0$. \monades{6}
\end{alist}